\chapter{Telecommunications: The Network War, in Lesser Detail}
\label{ch:telecom}

The three preceding chapters each took an industry from Western invention to Chinese supply-chain capture and priced every step. This chapter is deliberately shorter, and it is different in kind. Solar modules, battery cells and electric vehicles are \emph{commodities}: the campaign in each case commoditized the artifact and then owned what the artifact had to be made from. Telecommunications equipment is a \emph{network}---sold to a few hundred operators rather than millions of buyers, installed for decades, serviced continuously, and bound to standards that are written in committee rather than set by price. The playbook nevertheless ran here, one company executed almost all of it, and the record is complete enough to score. It belongs in Part~I for two reasons. It is the precedent closest in structure to what an ``AI stack'' export would be---infrastructure, financed, standardized, locked in by service---and it is the case most often invoked, in 2026, to describe what China is now supposed to be doing with sovereign AI. Part~III tests that invocation against what has actually been signed. This chapter supplies the template it is being tested against, and states at the end where the template does not transfer.

\section{The rise, measured}

Huawei was founded in 1987 as a reseller of imported switches and had about a fifth of the Chinese switch market by 1996. By 2006 revenue was US\$8.5B, by 2009 US\$21.8B, and by 2011 the company was at parity with Ericsson, the incumbent it had been imitating fifteen years earlier~\cite{ahrens-huawei} [A]. The global share series is the one that matters, and it is worth printing whole because its shape is the argument: 27.7\% of world telecom-equipment revenue in 2018, 28.1\% in the first half of 2019, 29\% in the second quarter of 2019 against Nokia's 15.7\% and Ericsson's 13.1\%, and 31\% in the first half of 2020---the year the sanctions bit hardest~\cite{delloro-huawei} [A]. Revenue peaked at RMB~891B in 2020, fell to RMB~636B in 2021 after the Honor divestiture and the loss of the handset business's chip supply, and recovered to RMB~880.9B (about US\$126B) in 2025 with net profit of RMB~68B and R\&D at 21.8\% of revenue~\cite{huawei-ar-2025} [P]. The company operates in more than 170 countries and regions by its own count~\cite{huawei-ar-2025}.

\section{The entry pattern: encircle the cities from the countryside}

The sequence is the playbook's target-selection and demand-before-supply moves executed against a market whose incumbents had chosen, rationally, not to serve its margins. Huawei went first where Ericsson and Alcatel would not go: the Chinese countryside, with more than two hundred staff across the counties of a single northeastern province against three or four for the incumbent; then Russia, Thailand, Brazil and South Africa; Ethiopia in 2003; a US\$200M CDMA network in Nigeria in 2005; the Netherlands and Germany by the middle of the decade; and only then the Tier-1 European operators, with British approval in 2005 as the pivot into the markets that set reference prices for everyone else~\cite{ahrens-huawei,dou-huawei} [A]. The Nortel episode belongs here as a caution rather than a proof: a 2004 intrusion traced to Shanghai removed some 1{,}488 documents including optical and next-generation wireless research, and Huawei underbid Nortel by about 40\% on a 2008 Canadian contract; but a 2014 University of Ottawa study found no evidence that the intrusion caused Nortel's 2009 bankruptcy, and this book does not assert what the record does not~\cite{dou-huawei} [A].

\section{The weapon: credit}

The instrument that distinguishes this campaign from solar and batteries, and that AI has not yet reproduced, is vendor finance at sovereign scale. China Development Bank opened a US\$10B credit facility for Huawei's customers in 2004 and raised it to \textbf{US\$30B in December 2009}, with a further US\$600M from the Export--Import Bank; roughly US\$10B had been drawn by 2013, and a 2011 European Union inquiry put the facility at about ten times Huawei's operating cash flow~\cite{ahrens-huawei,wsj-huawei-state} [A]. The \emph{Wall Street Journal}'s reconstruction in 2019 put total state support at up to US\$75B---some US\$46B in loans, credit lines and other financing from state lenders, about US\$25B in tax savings over 2008--18, plus grants and discounted land~\cite{wsj-huawei-state} [A]. Huawei disputed the characterization, and the counter-data must be printed with it: the same reporting has Huawei's prices about 30\% below rivals' in developing markets, but a 2017 European study found them only slightly lower in Europe and a 2019 study found them about 25\% \emph{higher} at home~\cite{wsj-huawei-state}. The discount was where the credit was. An operator in Lagos or Jakarta was not choosing between two switches; it was choosing between a switch it could finance and one it could not, and the financing came from the vendor's state.

This is the demand-before-supply move of Chapter~\ref{ch:playbook} run with a bank rather than a feed-in tariff, and it is the most exportable instrument in the playbook because it requires nothing of the recipient except a signature. It is also the instrument whose absence, in the AI campaign as of \datafreeze, Chapter~\ref{ch:nonaligned} records.

\section{The endgame: standards from the inside}

The consolidation-and-gatekeeping step took a form the earlier campaigns did not need. Rather than build a parallel standards body, Huawei captured position inside the existing one: first in declared 5G patent families and first in 3GPP technical contributions across 2015--24, with Chinese-headquartered firms above 40\% of declared 5G families~\cite{huawei-5g-seps} [A]. A vendor with the largest installed base \emph{and} the largest share of the specification is not a supplier the operator can leave at the next generation. That is move~M8 of the playbook executed in the only arena where it counts for a network, and it is the move the AI campaign has attempted and not yet achieved---Chapter~\ref{ch:harness} records Western agent-standards bodies with no Chinese member, and Chapter~\ref{ch:compute} an open instruction set with no ratified matrix extension.

\section{The pushback, and what it removed}

The bans came in sequence: Australia in August 2018, the US Entity List on 16 May 2019, the EU's 5G toolbox in January 2020, the United Kingdom's reversal on 14 July 2020, Sweden that October and, in January 2026, a Commission proposal for mandatory phase-out of high-risk suppliers~\cite{huawei-bans} [P/A]. The cost was real and mostly Western: the British removal was priced at up to \pounds2B and up to three years' delay; the American ``rip and replace'' programme had about US\$5B appropriated and 53 of 126 projects complete by June 2026 against a deadline it had already missed~\cite{huawei-bans} [A].

What the bans removed is the finding. Huawei's global share was 31\% in the first half of 2020 and \textbf{31\% again in the first half of 2025}---21\% outside China and 40\% outside North America~\cite{delloro-huawei} [A]. Chinese vendors' share of European 5G sites was about 32\% in mid-2022 and about 32\% at the end of 2024, with Germany at 59\%; of twenty-three European states that legislated, seven imposed restrictions~\cite{huawei-bans} [A]. Omdia counts Huawei the leading vendor in three of five world regions, third in Europe, and absent from North America~\cite{delloro-huawei}. The campaign failed where the security agencies had a veto and the operators could swap vendors without cost---the Five Eyes, the Nordics, Japan---and nowhere else. And the instrument that actually turned Britain was not a network-security finding but the May 2020 foundry rule: a semiconductor chokepoint, which is to say the Western counter-pattern that Chapter~\ref{ch:playbook} calls chokepoint denial, applied at the one layer where it still held~\cite{huawei-bans}.

\section{What transfers to AI, and what does not}

Scored against the six steps, telecommunications is a complete campaign: target selection (the unserved periphery), demand before supply (sovereign credit), overcapacity as strategy (the 30\% discount), selection (Huawei over ZTE, over a dozen domestic rivals), commoditization of the incumbents' rent (equipment margins compressed to the point that Nortel died and Alcatel and Lucent merged to survive), and an endgame of standards position and installed base that the bans could dent but not reverse. On the success--failure model of Chapter~\ref{ch:analogy} it sits in the kill zone on modularity and codification and \emph{outside} it on trust and certification, which is exactly the partition the outcome shows: won everywhere trust was not a procurement criterion, lost everywhere it was.

Three things transfer directly to the AI campaign, and Parts~II and~III document each. The rural-first entry pattern is the non-aligned adoption channel of Chapter~\ref{ch:nonaligned}: usage arrives before procurement, in markets the incumbents price for someone else. The standards-from-inside endgame is the move the harness and compute chapters record as attempted. And the demonstration that a chokepoint at the silicon layer can reverse a network position built over two decades is the strongest single argument, in this book, for taking the semiconductor chapter's constraint seriously.

Three things do not transfer, and Chapter~\ref{ch:nonaligned} tests each against the record at \datafreeze. \emph{The credit is absent.} Chapter~\ref{ch:nonaligned} finds no sovereign credit facility for AI infrastructure exports in the record, and the one contracted deployment on Chinese silicon abroad financed by the recipient. \emph{The capacity is absent.} Huawei exported switches because it had surplus factories and a state that wanted them running; in AI the same company is supply-constrained at the accelerator, a single domestic customer's order absorbs a year of output, and Beijing's own mandates pull the chips home. \emph{And the artifact does not lock in.} A telecom network is hardware plus service plus debt, and it binds the operator for a generation. Open weights are the opposite object: they cost the recipient nothing and bind the recipient to nothing, and they are the instrument through which Chinese AI actually diffuses---one hundred and fifty thousand derivatives of a single model family, five times the download volume of its nearest Western rival~\cite{qwen-derivatives} [V]. That is a more effective diffusion mechanism than Huawei ever had and a far weaker rent-capture mechanism, and the difference between the two is the reason this chapter is shorter than the three before it. The telecom war shows what the playbook does with a network. The AI war, so far, is being fought with a commons.
